% system-model.tex -- Enterprise Node Orchestrator Paper (RU-V5)

\section{System Model}

\subsection{Network Architecture}

The Basis Network enterprise ZK validium operates on a three-tier architecture:

\begin{enumerate}
\item \textbf{Enterprise Systems Layer}: Client applications (PLASMA for
  industrial maintenance, Trace for SME ERP) that produce business
  transactions. These systems hash all data client-side before submission;
  the node never receives raw enterprise data.

\item \textbf{Validium Node Layer}: A single Node.js service per enterprise
  that maintains off-chain state (Sparse Merkle Tree), aggregates
  transactions into batches, generates Groth16 ZK proofs, distributes
  witness data to a Data Availability Committee (DAC), and submits proofs
  to the settlement layer.

\item \textbf{Settlement Layer}: The Basis Network L1, an Avalanche
  Subnet-EVM chain (Chain ID 43199) with zero-fee permissioned transactions
  and sub-second Snowman consensus finality.
\end{enumerate}

\subsection{Trust Assumptions}

\begin{itemize}
\item The enterprise node operator is trusted for liveness but not for
  data integrity (proofs enforce correctness).
\item The DAC operates under a $t$-of-$n$ threshold (default 2-of-3):
  data remains available if at least $t$ committee members are honest.
\item The L1 settlement layer is assumed correct (Avalanche consensus
  with permissioned validator set).
\item Enterprise client systems are trusted to correctly hash data
  before submission (the node verifies signatures but not data content).
\end{itemize}

\subsection{Privacy Model}

The zero data leakage invariant requires that no raw enterprise data
crosses the node boundary. The data flow is:

\begin{enumerate}
\item Enterprise systems compute $h = \text{SHA-256}(\text{record data})$ and
  submit transaction tuples $(k, h, \sigma)$ where $k$ is a deterministic
  key derived from the record identifier and $\sigma$ is the enterprise
  signature.

\item The node inserts $(k, h)$ into the Sparse Merkle Tree, generating
  state root transitions $r_{\text{prev}} \to r_{\text{new}}$.

\item The ZK proof reveals only: $r_{\text{prev}}$, $r_{\text{new}}$,
  batch number, and enterprise identifier. No transaction keys, values,
  or counts are disclosed.

\item Witness data is distributed to the DAC via Shamir secret sharing
  (information-theoretically secure: no single committee member can
  reconstruct the witness).

\item The L1 contract stores only: state roots, proofs, and batch metadata.
\end{enumerate}

\subsection{Performance Requirements}

The orchestrator must satisfy:
\begin{itemize}
\item \textbf{End-to-end latency}: $< 90$~s for a batch of 64 transactions
  (60~s proving budget + 30~s overhead budget).
\item \textbf{Continuous availability}: The node must accept new transactions
  during proof generation and L1 submission (no downtime windows).
\item \textbf{Crash recovery}: Zero transaction loss under arbitrary crash
  scenarios, with bounded recovery time.
\item \textbf{Memory footprint}: $< 500$~MB for a single enterprise node
  with up to 100K state entries.
\end{itemize}
